В работе рассмотрена задача автоматизированной оцифровки химической литературы, включающая обработку сканированных страниц, выделение структурных областей документа, распознавание текстовой информации и извлечение химических структур в машиночитаемые форматы. Показано, что специфика химических изданий связана не только с наличием обычного текста, но и с большим количеством структурных формул, реакционных схем, таблиц, подписей и графических элементов, что делает применение классических OCR-подходов недостаточным.

Предложенный подход основан на последовательной обработке документа: предварительном улучшении качества изображений, сегментации страниц, распознавании текстовых блоков, обнаружении химических структур и их последующем преобразовании в стандартизированный формат представления SMILES. Такой конвейер позволяет перейти от неструктурированного изображения страницы к набору данных, пригодных для дальнейшего поиска, анализа, индексации и использования в информационных системах.
Особое внимание уделено применению методов машинного зрения и современных моделей распознавания химических структур. Использование OCR-средств для текста и OCSR-моделей для химической графики позволяет разделить обработку различных типов информации и повысить качество итогового извлечения данных. Дополнительная постобработка с использованием эвристик и языковых моделей играет важную роль при исправлении ошибок распознавания, нормализации результатов и повышении согласованности извлеченной информации.

В результате проведенного исследования показана практическая значимость комплексного подхода к оцифровке химической литературы. Разработанная схема обработки может быть использована как основа для создания программной системы, обеспечивающей автоматическое извлечение текстовой и структурной химической информации из сканированных документов. Перспективными направлениями дальнейшей работы являются расширение набора тестовых данных, повышение точности сегментации сложных страниц, улучшение качества распознавания дефектных изображений, а также интеграция полученных данных в системы семантического поиска и анализа химической информации.